\section{Related Work and Positioning}
\label{sec:related_work}

\subsection{PAC-Bayes under Dependence and Sequential Observation}
PAC-Bayes for dependent observations is established. Chromatic and blocking arguments cover dependency graphs and stationary mixing processes \cite{ralaivola2010chromatic}; oracle inequalities have been developed for weakly dependent time-series prediction \cite{alquier2012model,alquier2013prediction}; and recent results address Markov models and stable nonlinear dynamical systems \cite{banerjee2021tempered,eringis2024dynamical}. These approaches typically require explicit dependence, stationarity, boundedness, or stability assumptions.

Martingale and supermartingale methods provide a different route. Seldin \emph{et al.} derived PAC-Bayesian inequalities for martingales \cite{seldin2012martingales}. Online and time-uniform constructions subsequently allowed data-dependent posterior sequences while keeping the prior predictable \cite{haddouche2022online,chugg2023unified}. We use this route because overlapping lag windows form an adapted sequence, and because it avoids treating an estimated mixing coefficient as known. The bounded-loss formula used here is a direct specialization of this literature, not a new general concentration inequality.

\subsection{Computable and Non-Vacuous Certificates}
A formal inequality becomes useful only when the prior, posterior, and empirical term can be computed without hidden data reuse. Data-dependent priors can be valid when their information source is explicitly separated or otherwise controlled \cite{parrado2012dataprior,rivasplata2020beyond}. Work on neural-network certificates shows that prior localization and posterior-scale selection are decisive for tightness \cite{perezortiz2021tighter}. Foong \emph{et al.} show that small-sample PAC-Bayes bounds may remain looser than direct test bounds even after careful optimization \cite{foong2021small}. Heavy-tail and unbounded-loss PAC-Bayes results offer alternatives to clipping \cite{haddouche2021unleashed,haddouche2023heavy,rodriguez2024more}; the current implementation does not claim those guarantees and instead reports clipping explicitly.

\subsection{Sparse TSK Systems}
Bayesian TSK models, block-sparse dynamic fuzzy models, and sparse rule or coefficient learning are well developed \cite{gu2018btsk,li2018blocksparse,li2022bayesiants,lughofer2010sparsefis,luo2013hierarchical,wiktorowicz2022timeseries}. Clustering-based rule construction and online TSK identification are also classical \cite{Chiu1994Cluster,angelov2004online}. Existing work generally targets prediction, parameter estimation, or interpretability; it does not provide the specific audited chain from a prior-only radius cover to a realized randomized dimension, KL cost, and sequential certificate.

\begin{table*}[t]
\centering
\caption{Positioning relative to representative PAC-Bayes literature. ``This work'' uses a known time-uniform supermartingale route; its contribution is the chronological, computable specialization and the structural/operational evaluation.}
\label{tab:prior_work}
\small
\setlength{\tabcolsep}{4pt}
\begin{tabular}{p{2.5cm}p{2.2cm}p{2.4cm}p{3.0cm}p{5.4cm}}
\toprule
Work & Data regime & Loss / guarantee & Main computational focus & Difference from the present study \\
\midrule
Seldin et al. \cite{seldin2012martingales} & Martingales & PAC-Bayes martingale inequalities & General concentration & Provides foundational martingale PAC-Bayes inequalities; no chronological TSK construction or rule-count audit. \\
Chugg et al. \cite{chugg2023unified} & Sequential, non-IID & Time-uniform PAC-Bayes recipes & General supermartingale framework & The bound used here is a bounded-loss specialization, not a competing general theorem. \\
P\'erez-Ortiz et al. \cite{perezortiz2021tighter} & IID classification & Tight neural-network risk certificates & Joint prior/posterior optimization & Motivates computable non-vacuous certificates, but not temporal dependence or fuzzy structure. \\
Foong et al. \cite{foong2021small} & IID, small samples & Tightness limits & Bound-vs-test comparison & Motivates reporting vacuity and clipping; it does not address sequential dependence. \\
Haddouche--Guedj \cite{haddouche2023heavy} & Sequential/heavy-tailed & Supermartingale bounds for heavy tails & Unbounded-loss control & Offers a route beyond clipping; the current implementation deliberately uses bounded clipped loss. \\
Alquier--Wintenberger \cite{alquier2012model} and Eringis et al. \cite{eringis2024dynamical} & Weakly dependent time series / stable systems & Generalization under dependence & Rates or stability/mixing conditions & These require dependence or stability assumptions; the primary certificate here uses predictability and bounded loss instead. \\
This work & Overlapping, possibly nonstationary forecast windows & Conditional sequential clipped risk & Hierarchical structural coding; radius-vs-fixed-$K$ ablation; frozen deployment gate & Empirically separates structural sparsity from activation sparsity and tests a negative development gate followed by an independent confirmatory energy case. \\
\bottomrule
\end{tabular}
\end{table*}


\subsection{Precise Novelty Claim}
The present study does not introduce a new clustering algorithm, a new general PAC-Bayes theorem, or a proof that TSK models dominate linear or seasonal forecasts. Its contribution is the intersection of: (i) a chronologically valid finite prior over complete forecasting structures; (ii) a computable sequential certificate for a linear-in-consequents TSK model; (iii) a controlled fixed-$K$ versus radius-controlled structural ablation; and (iv) an operational protocol that records a negative development case before an independent frozen confirmatory case. Table~\ref{tab:prior_work} makes the distinction from the closest PAC-Bayes literature explicit.
